
\documentclass{report}

\input{preamble}
\input{macros}
\input{letterfonts}

\title{\Huge{Analisi Matematica}}
\author{\huge{Caberlotto Giovanni}}
\date{}

\begin{document}

\maketitle
\newpage% or \cleardoublepage
% \pdfbookmark[<level>]{<title>}{<dest>}
\pdfbookmark[section]{\contentsname}{toc}
\tableofcontents
\pagebreak

\chapter{Teoria degli Insiemi ed Elementi di Logica} % (fold)
\label{chap:Teoria degli Insiemi ed Elementi di Logica}

\section{Definizione di Insieme}

\section{Rappresentazione di un Insieme}

\section{Tipi di Insiemi}

\section{Operazioni tra Insiemi}

\section{Numeri Reali}

\section{Intervalli ed Estremi}

\section{Intorno di un punto}



% chapter Teoria degli Insiemi ed Elementi di Logica (end)

\chapter{Funzioni Elementari Grafici e Proprietà} % (fold)
\label{chap:Funzioni Elementari Grafici e Proprietà}

\section{Radicali}

\section{Logaritmo}

\section{Esponenziali}

\section{Valore Assoluto}

\section{Seno e Coseno}

\section{Tangente e Cotangente}

\section{Funzione Costante}

\section{Funzione Identità}

\section{Funzione Lineare}

\section{Funzione Polinomiale}

% chapter Funzioni elementari proprietà e analisi grafico (end)

\chapter{Teoria delle Funzioni} % (fold)
\label{chap:Teoria delle Funzioni}

% chapter Teoria delle Funzioni (end)

\chapter{Limiti} % (fold)
\label{chap:Limiti}

% chapter Limiti (end)

\chapter{Derivate} % (fold)
\label{chap:}

% chapter  (end)

\chapter{Studio di funzione} % (fold)
\label{chap:Studio di funzione}

% chapter Studio di funzione (end)

\chapter{Integrali} % (fold)
\label{chap:Integrali}

% chapter Integrali (end)

\chapter{Successioni} % (fold)
\label{chap:Successioni}

% chapter Successioni (end)

\chapter{Serie} % (fold)
\label{chap:Serie}

% chapter Serie (end)

\chapter{Numeri Complessi} % (fold)
\label{chap:Numeri Complessi}

% chapter Numeri Complessi (end)

\end{document}
